%=====================================================================
% 光学（02）Lecture 3 · 作业答案卡（English，供手抄上交）
% 来源：Polarization Devices 1.pdf p110（三个偏振器件与干涉问题）
% 内容与 L3 讲义附录 D 一致（外教L3-偏振器件与干涉-精讲讲义-20261116.pdf p30--p31）
% 编译：xelatex 两遍
%=====================================================================
\documentclass[10.5pt]{ctexart}
%=====================================================================
% theme-violet.tex —— 讲义精讲 · 紫花主题
% 用法：\documentclass[11pt]{ctexart}
%       \input{assets/theme-violet}
%       \begin{document} ... \end{document}
% 只用 ctexart 自带能力 + xcolor/enumitem/fancyhdr/pifont/amsmath，
% 不依赖 tikz/tcolorbox，保证 MiKTeX 可直接编译。
%=====================================================================

\usepackage[a4paper,left=2.0cm,right=2.0cm,top=2.1cm,bottom=2.3cm,%
            headsep=12pt,footskip=22pt]{geometry}
\usepackage{amsmath,amssymb}
\usepackage{needspace}
\usepackage{xcolor}
\usepackage{array}
\usepackage{enumitem}
\usepackage{pifont}
\usepackage{fancyhdr}
\usepackage[hidelinks]{hyperref}

%---------------------------------------------------------------- 配色
\definecolor{violet}{RGB}{107,78,143}      % 主色：紫
\definecolor{violetdeep}{RGB}{63,44,92}    % 深紫：标题/强调
\definecolor{petal}{RGB}{183,159,208}      % 花瓣：装饰线、花饰
\definecolor{bandviolet}{RGB}{247,243,251} % 浅紫底：档案/提示框
\definecolor{bandcream}{RGB}{252,249,242}  % 米白底：易错/注意框
\definecolor{rulepurple}{RGB}{217,207,230} % 分隔线
\definecolor{inkgray}{RGB}{92,96,110}      % 次要文字

%---------------------------------------------------------------- 字体
% 有 Noto CJK 用 Noto，没有就退回 ctex 默认（Windows 下为宋体/黑体）
\IfFontExistsTF{Noto Serif CJK SC}{\setCJKmainfont{Noto Serif CJK SC}}{}
\IfFontExistsTF{Noto Sans CJK SC}{\setCJKsansfont{Noto Sans CJK SC}}{}

%---------------------------------------------------------------- 版式
\setlist{nosep,leftmargin=1.9em}
\linespread{1.28}
\setlength{\parskip}{3.5pt}
\setlength{\parindent}{0pt}
\setlength{\tabcolsep}{5pt}
\renewcommand{\arraystretch}{1.28}

%---------------------------------------------------------------- 花饰
% \ding{95} = 五瓣小花，\ding{96} = 星形花；都来自 pifont，必定存在
\newcommand{\orn}{\textcolor{petal}{\ding{95}}}          % 行内小标
\newcommand{\ornstar}{\textcolor{petal}{\ding{96}}}      % 星形花
\newcommand{\flowerrule}{%                                % 居中花饰分隔线
  \par\vspace{7pt}%
  \noindent\makebox[\linewidth]{\color{petal}%
    \rule{0.20\linewidth}{0.6pt}\hspace{8pt}{\small\ding{95}}\hspace{4pt}%
    \raisebox{1.1pt}{\tiny\ding{95}}\hspace{4pt}{\small\ding{95}}\hspace{8pt}%
    \rule{0.20\linewidth}{0.6pt}}%
  \par\vspace{7pt}}

%---------------------------------------------------------------- 标题与节头
\newcommand{\doctitle}[3]{%   #1 主标题  #2 副标题  #3 元信息（可用 \\ 分行）
  \begin{center}
    {\color{petal}\small \orn\enspace #2\enspace \ornstar}\\[7pt]
    {\LARGE\bfseries\color{violetdeep} #1}\\[10pt]
    {\color{petal}\rule{0.55\linewidth}{1pt}}\\[4pt]
    {\small\color{inkgray} #3}
  \end{center}\vspace{2pt}}

\newcommand{\hdr}[2]{%   #1 章节号（如 §0）  #2 标题
  \par\vspace{12pt}%
  \noindent\colorbox{violetdeep}{\color{white}\bfseries\small\;#1\;}\hspace{7pt}%
  {\color{violetdeep}\large\bfseries #2}\par\nobreak
  \vspace{3pt}{\color{petal}\hrule height 1.2pt}\vspace{6pt}}

\newcommand{\subhdr}[1]{%
  \par\vspace{6pt}%
  \noindent{\color{violet}\bfseries #1}\par\vspace{1pt}}

%---------------------------------------------------------------- 提示框
\newcommand{\callout}[1]{% 浅紫底：讲义档案、读法说明
  \par\medskip\noindent
  \fcolorbox{rulepurple}{bandviolet}{\parbox{\dimexpr\linewidth-2\fboxsep-2\fboxrule}{\vspace{2pt}#1\vspace{2pt}}}%
  \par\medskip}

\newcommand{\warnbox}[1]{% 米白底：易错、边界条件
  \par\medskip\noindent
  \fcolorbox{rulepurple}{bandcream}{\parbox{\dimexpr\linewidth-2\fboxsep-2\fboxrule}{\vspace{2pt}#1\vspace{2pt}}}%
  \par\medskip}

\newcommand{\gapitem}[2]{% #1 缺口类型（紫） #2 内容
  \par\vspace{2pt}\noindent
  \hangindent=1.4em\hangafter=1
  {\color{petal}\small\ding{96}}\enspace{\bfseries\color{violetdeep}#1}\enspace #2\par}

% 题干块（\qstem）：每道例题／作业的「详解」开头必须先复述题干与出处。
% 起因：2026-10-25 第五章讲义把七道作业的题面集中放在一页、详解另起一页，
% 读者翻到「第 4 题」时看不到条件，等于一道没有题干的题目。
% 规则：详解块内自带题干，读者不翻页就能读懂题在问什么。
\newcommand{\qstem}[1]{%
  \par\medskip\noindent
  \fcolorbox{violet}{bandviolet}{%
    \parbox{\dimexpr\linewidth-2\fboxsep-2\fboxrule}{%
      \vspace{2pt}%
      {\small\bfseries\color{violet}题干}\hspace{0.8em}#1%
      \vspace{2pt}}}%
  \par\medskip}

% v2: 保护题干与开头，长题可跨页，续页补上下文。
\newcommand{\qbreak}{\needspace{6\baselineskip}}
% 独立练习：ID、来源、题干、考点、答案标签（通常 ans:ID）。
\newcommand{\practice}[5]{%
  \qbreak\subhdr{练习 #1}\qstem{（#2）#3}%
  \textbf{考点}：#4\par
  详解见附录第~\pageref{#5}~页。\par\medskip}
% 附录开头：ID、来源、条件摘要；自动产生 ans:ID 标签。
\newcommand{\answerstart}[3]{%
  \qbreak\subhdr{#1\quad 详解}\label{ans:#1}\qstem{（#2）#3}}
\newcommand{\answercontinue}[2]{%
  \needspace{5\baselineskip}\subhdr{#1\quad 续解}%
  {\small\color{inkgray}本页接续：#2}\par\medskip}
% 延迟提供默认值，让旧讲义可在导言区保留自己的 \newcommand{\step}。
\AtBeginDocument{\providecommand{\step}[1]{\needspace{3\baselineskip}\par\medskip\textbf{#1}\par}}

% 前提旁注：只放在「需要理解的那一步」后面，给陈述 + 本题代入，不重证。
\newcommand{\aside}[1]{%
  \par\vspace{1pt}\noindent
  \hangindent=1.6em\hangafter=1
  {\small\color{inkgray}%
    \hspace*{0.6em}{\color{petal}$\triangleright$}\hspace{2pt}#1}%
  \par\vspace{2pt}}

%---------------------------------------------------------------- 表格
% 本机 MiKTeX 的 colortbl 与 array 版本不匹配（\rowcolor 会报错），
% 故表格统一用「上粗紫线 + 表头 + 中细线」的三线式，不依赖 colortbl。
\newcommand{\tabtop}{\par\vspace{3pt}{\color{violet}\hrule height 1.2pt}\vspace{1pt}}
\newcommand{\tabmid}{\par\vspace{1pt}{\color{petal}\hrule height 0.5pt}\vspace{1pt}}
\newcommand{\tabbot}{\par\vspace{3pt}{\color{rulepurple}\hrule height 0.5pt}\vspace{8pt}}
\newcommand{\thead}[1]{{\bfseries\color{violetdeep}#1}}

%---------------------------------------------------------------- 页眉页脚
\pagestyle{fancy}
\fancyhf{}
\renewcommand{\headrulewidth}{0.4pt}
\renewcommand{\headrule}{\hbox to\headwidth{\color{rulepurple}\leaders\hrule height \headrulewidth\hfill}}
\fancyhead[L]{\small\color{inkgray}\songti \doctag}
\fancyhead[R]{\small\color{inkgray}\songti \coursetag}
\fancyfoot[C]{\small\color{inkgray}\orn\enspace{\color{violet}\thepage}\enspace\orn}

% 每篇讲义在导言区设这两个值
\newcommand{\doctag}{}
\newcommand{\coursetag}{}
\newcommand{\settags}[2]{\renewcommand{\doctag}{#1}\renewcommand{\coursetag}{#2}}

%---------------------------------------------------------------- 数学记号
\newcommand{\vect}[1]{\boldsymbol{#1}}
\newcommand{\dd}{\,\mathrm{d}}
\newcommand{\ddx}{\frac{\mathrm{d}}{\mathrm{d}x}}
\addtolength{\jot}{2pt}



\setlength{\headheight}{20pt}
\renewcommand{\tabmid}{\noalign{\vskip 2pt{\color{petal}\hrule height 0.5pt}\vskip 3pt}}
\settags{偏振器件 · L3 作业答案卡}{光学（02）· 外教 Lecture 3}
\geometry{left=1.9cm,right=1.9cm,top=1.8cm,bottom=2.0cm}
\setlength{\parskip}{2.2pt}
\begin{document}

\begin{center}
{\LARGE\bfseries\color{violetdeep} Homework: Polarization Devices \& Interference}\\[5pt]
{\color{petal}\small \ding{95}\enspace Answer sheet (courseware p110)\enspace \ding{96}}\\[4pt]
{\small\color{inkgray}Handwrite in English;\ keep the four parts \textbf{Given} / \textbf{Formulas used} /
\textbf{calculation} / \textbf{Answers with units};\ sign your name in English.}
\end{center}
\vspace{-2pt}
{\color{petal}\hrule height 1pt}

\subhdr{Problem 1 \ ---\ two (or three) polarizing prisms}

\textbf{Given}\quad Unpolarized light of intensity $I_0$; transmission planes of the two prisms at
$\theta=45^{\circ}$; in (b) the second prism is rotated to $60^{\circ}$; in (c) a third prism is inserted at
$30^{\circ}$ to the first.

\textbf{Formulas used}\quad Unpolarized light through the first polarizer: $I_1=I_0/2$;\quad
linearly polarized light through an analyzer (Malus' law): $I=I_{\text{before}}\cos^{2}\theta$.

\textbf{Calculation}\quad
(a) The first prism transmits one half; the second follows Malus' law:
\begin{equation*}
I_2=\frac{I_0}{2}\cos^{2}45^{\circ}=\frac{I_0}{2}\times\frac12=\frac{I_0}{4}=0.25\,I_0 .
\end{equation*}
(b) Only the angle changes:
\begin{equation*}
I_2=\frac{I_0}{2}\cos^{2}60^{\circ}=\frac{I_0}{2}\times\frac14=\frac{I_0}{8}=0.125\,I_0 .
\end{equation*}
(c) With three prisms the angles must be taken \emph{between neighbouring prisms}: $30^{\circ}$ from the first
to the third, and $45^{\circ}-30^{\circ}=15^{\circ}$ from the third to the second:
\begin{equation*}
I_3=\frac{I_0}{2}\cos^{2}30^{\circ}\cos^{2}15^{\circ}
=\frac{I_0}{2}\times0.750\times0.933\approx0.350\,I_0 .
\end{equation*}

\textbf{Explanation}\quad Each polarizer projects the vibration onto its own transmission axis. The inserted
prism turns the vibration to $30^{\circ}$, so the projections onto the second prism no longer vanish; neither
$\cos30^{\circ}$ nor $\cos15^{\circ}$ is zero. Intensities in a chain of polarizers are therefore \emph{not}
additive, and a prism that would block light on its own lets some light through when it is inserted between
two crossed ones.

\textbf{Answers}\quad (a) $\dfrac{1}{4}I_0$ (25\%);\qquad (b) $\dfrac{1}{8}I_0$ (12.5\%);\qquad
(c) $\approx0.35\,I_0$ (35\%), because the vibration is re-projected at each prism.

\subhdr{Problem 2 \ ---\ half-wave plate of quartz}

\textbf{Given}\quad $\lambda=633\ \mathrm{nm}$, quartz with $\Delta n=0.0091$, half-wave plate
($\delta=\pi$); in (b) the incident azimuth is $20^{\circ}$ to the fast axis.

\textbf{Formulas used}\quad Retardation $\Delta=\Delta n\,d$ and phase difference
$\delta=\dfrac{2\pi}{\lambda}\Delta n\,d$;\quad half-wave condition $\Delta n\,d=\dfrac{\lambda}{2}$;\quad
azimuth after a half-wave plate $\alpha_{\text{out}}=2\theta-\alpha_{\text{in}}$, where $\theta$ is the angle
of the fast axis.

\textbf{Calculation}\quad
(a)
\begin{equation*}
d=\frac{\lambda}{2\Delta n}=\frac{633\ \mathrm{nm}}{2\times0.0091}
=\frac{633}{0.0182}\ \mathrm{nm}\approx34\,780\ \mathrm{nm}\approx34.8\ \mu\mathrm{m}
\end{equation*}
which is the zero-order ($m=0$) plate, i.e.\ the thinnest possible one.
(b) Take the fast axis as $0^{\circ}$, so $\theta=0^{\circ}$ and $\alpha_{\text{in}}=20^{\circ}$:
\begin{equation*}
\alpha_{\text{out}}=2\times0^{\circ}-20^{\circ}=-20^{\circ}.
\end{equation*}
The transmitted light is still linear, with its azimuth mirrored about the fast axis.
(c) A half-wave plate reverses the sense of rotation, so circularly polarized light leaves circularly polarized
with the \emph{opposite handedness} (left $\leftrightarrow$ right).
(d) The same plate at $450\ \mathrm{nm}$ no longer gives half a wave:
\begin{equation*}
\frac{\Delta}{\lambda}=\frac{0.0091\times34\,780\ \mathrm{nm}}{450\ \mathrm{nm}}
=\frac{316.5}{450}\approx0.703
\quad\Longrightarrow\quad
\delta=0.703\times360^{\circ}\approx253^{\circ},
\end{equation*}
which differs from $180^{\circ}$ by $73^{\circ}$: the emergent light is elliptical, not linear, so the plate
does not act as a half-wave plate at this wavelength.

\textbf{Answers}\quad (a) $d_{\min}\approx34.8\ \mu\mathrm{m}$;\qquad
(b) $\alpha_{\text{out}}=-20^{\circ}$ (mirror image about the fast axis);\qquad
(c) circular polarization of the opposite handedness;\qquad
(d) at $450\ \mathrm{nm}$ the retardation is $0.70\lambda$ ($\delta\approx253^{\circ}$), so the output is elliptical.

\subhdr{Problem 3 \ ---\ birefringent plate between crossed polarizers}

\textbf{Given}\quad $d=30\ \mu\mathrm{m}=3.0\times10^{4}\ \mathrm{nm}$, $\Delta n=0.0091$, crossed polarizers,
fast axis at $45^{\circ}$ to the polarizer, white light ($400$--$700\ \mathrm{nm}$).

\textbf{Formulas used}\quad $\Delta=\Delta n\,d$;\quad transmitted intensity between crossed polarizers
$I=I_0\sin^{2}2\varphi\,\sin^{2}\dfrac{\delta}{2}$ with $\delta=\dfrac{2\pi}{\lambda}\Delta$;\quad
extinction when $\Delta=m\lambda$, i.e.\ $\lambda_m=\Delta/m$.

\textbf{Calculation}\quad
(a)
\begin{equation*}
\Delta=\Delta n\,d=0.0091\times30\ \mu\mathrm{m}=0.273\ \mu\mathrm{m}=273\ \mathrm{nm}.
\end{equation*}
(b) Extinction requires $\Delta=m\lambda$ with integer $m$, so $\lambda_1=273\ \mathrm{nm}$,
$\lambda_2=136.5\ \mathrm{nm}$, $\lambda_3=91\ \mathrm{nm}$, \dots\ — all in the ultraviolet.
\textbf{No wavelength in $400$--$700\ \mathrm{nm}$ is extinguished.} (At $546\ \mathrm{nm}$,
$\Delta/\lambda=1/2$ and $\sin^{2}(\pi/2)=1$, i.e.\ a \emph{maximum}.)
(c) The transmitted light is pale yellow-green: $\sin^{2}(\pi\Delta/\lambda)$ reaches $1$ near
$546\ \mathrm{nm}$ and falls to about $0.89$ at $700\ \mathrm{nm}$ and $0.71$ at $400\ \mathrm{nm}$, so red and
blue are attenuated more than green. With monochromatic light of $\lambda=273\ \mathrm{nm}$ (the $m=1$
extinction, in the ultraviolet) the field is black because $\Delta=\lambda$, hence $\delta=2\pi$ and
$\sin^{2}(\delta/2)=0$.
(d) A wedge has a thickness that grows linearly along one direction, so $\Delta$ and $\delta$ also vary
linearly: the field shows a set of \textbf{straight fringes parallel to the wedge edge}. Each fringe is a line
of constant retardation (an isochromat) — bright and dark bands in monochromatic light, a sequence of coloured
bands in white light (Newton's scale).

\textbf{Answers}\quad (a) $\Delta=273\ \mathrm{nm}$;\qquad (b) none in the visible range (the extinctions lie
at $273\ \mathrm{nm}$ and below);\qquad (c) pale yellow-green, black at $273\ \mathrm{nm}$ because
$\delta=2\pi$;\qquad (d) straight fringes parallel to the wedge edge.

\vspace{4pt}
{\color{petal}\hrule height 0.6pt}
{\small\color{inkgray}\textbf{Before handing in:} handwritten; four parts per problem; answers with units;
name signed in English; placed on the instructor's desk before the lecture starts.}

\end{document}
